\section{Original Contribution}

The dissertation focuses on two types of delay that are peculiar to the \gls{rl} setting: state observation delay and action execution delay. 
Our goal is to broaden the understanding of such delays, as well as to provide efficient solutions. 
A particular focus is placed on constant delays that are ubiquitous in the literature and are a good approximation of the real delay in many cases.
However, we will repeatedly study the theoretical and empirical impact of extending the framework to more original types of delay.
The contributions of this thesis follow four objectives: providing a unifying framework for delays in \gls{rl}; giving a theoretical ground for better understanding delays; designing theoretically grounded and efficient algorithms to tackle delayed environments and providing an extensive empirical evaluation of the proposed algorithms, as well as algorithms from the literature.
In the following, the contributions for each of these objectives will be detailed before referring the reader to the publications that I have co-authored and where these contributions can be found.

\subsubsection{Unifying Framework for Delays}
In this dissertation, we provide an extensive review of the literature on delays, in \gls{rl} and related fields.
The focus is on state observation and action execution delay in \gls{rl} but other types of delay are also discussed.
Notably, we gather algorithms from the literature into three categories.
The first category is the \emph{augmented state} approach. 
The idea is to augment the last observed state with the sequence of actions that have been selected by the agent, but whose outcome has not been observed yet.
Using this new state allows one to cast the delayed problem back to the traditional undelayed \gls{mdp} framework.
The second category is the \emph{memoryless} approach, where the agent ignores the delay and uses only information about the last observed state.
The idea is simple, yet it can be efficient in practice.
Finally, the \emph{model-based} approach learns a model of the environment to predict the state of the environment at which the actions will apply.
To the best of our knowledge, there are no previous works providing an overview of delays in \gls{rl}. 
We believe that providing this information is beneficial for the field, as it is not uncommon to find overlapping results in the literature.

In addition, we propose a unified mathematical framework for the problem of delay.
This proposition follows the same line of thought that the delayed literature would benefit from more structure.
Our framework contains all the most common types of delay as a particular case and naturally interlocks with the \gls{mdp} model.

\subsubsection{Theoretical Understanding of Delays}
First, we analyse the theoretical effect of delayed state observation or action execution on the \gls{mdp}.
Our first result is a formal demonstration of a seemingly evident, yet left unproved, result: the longer the delay, the lower the maximum attainable cumulative reward--or return.
In the same vein, we show how the delay affects the variability of returns that the set of policies can achieve. 
A longer delay reduces the range of return that the agent can get, but also reduces the variance of its return against repeated learning of the same task. 
Then, we extend the constant delay framework to a multiple-action delay one, where a single action can impact several transitions in the future. 
This framework is interesting from a theoretical perspective, as it contains the constantly delayed process as a particular case. 
Moreover, it naturally extends the higher order \glspl{mc} model to \glspl{mdp}.
We show that multiple action-delayed processes can be cast back to an \gls{mdp} by state augmentation, similarly to constantly delayed processes. 
Then we study the properties and peculiarities of this framework with respect to the traditional \gls{mdp}. 
In particular, this allows us to underline the connection between \gls{mdp} of higher order and \gls{mdp} with delay.
Notably, some properties, such as being unichain, are shown to not transfer from the underlying \gls{mdp} to the delayed process built on it, while others do, such as communicating. 

Second, we shift the focus to delayed \gls{rl} algorithms and study their guarantees. 
We analyse the approaches to delays, which are the augmented state, the memoryless, and the model-based approaches.
It is known that the augmented state approach allows one to find the optimal policy for the delayed problem. 
Instead, we show that the model-based approach can be too restrictive and discard the optimal policy.
However, in the context of a smooth \gls{mdp}--concept that will be defined in the dissertation--we demonstrate that model-based approaches provide a nearly optimal policy.
We extend the previous guarantees to non-integer delays. 
Non-integer delays are reminiscent of the asynchronous environment hypothesis of real-time \gls{rl}. 
The agent sees the state at a given frequency, whereas it can only act with a phase change. 
Finally, we derive performance bounds for the case where the model-based policy has been trained on a constant delay while it is tested on stochastic delays.

\subsubsection{Algorithms for Delayed Reinforcement Learning}
On the algorithmic side, we first explore model-based approaches, where we design an algorithm capable of learning a vectorial representation of the probability distribution--or belief--of the current state.
To learn a representation of the belief, a first neural network, a Transformer, processes the augmented state and produces a representation for each future state up to the undelayed one.
This representation is trained to correspond to the belief by another network that fits the distribution, a normalising flow network. 
This representation can be plugged into an \gls{rl} algorithm to increase performance and boost learning speed compared to using the augmented state directly.
The effect is even clearer in stochastic environments.
Moreover, we provide a straightforward approach for utilising this algorithm, which has been pre-trained on a specific delay, to adapt to shorter delays.


Then we explore a different approach to learn a policy by imitating the behaviour of an undelayed expert. 
Practically, the delayed agent is trained to replicate the actions selected by the undelayed agent given the current state, but using delayed information only--the augmented state.
We show how to take advantage of the theoretical guarantees obtained for smooth \glspl{mdp} together with the theoretical results from imitation learning to guide the design of the algorithm.
We show empirically that this approach yields state-of-the-art performance and is versatile enough to apply to constant, non-integer and stochastic delays. 

As a third approach, we explore memoryless approaches.
In this case, the problem becomes inherently non-stationary. 
Therefore, we design a theory on how to adapt to future non-stationarity by optimising a probabilistic lower bound of the estimation of future performance.
This estimation of the future performance is based on multiple importance sampling to account for the variety of past dynamics due to non-stationarity.
Under smooth non-stationary, the bias of this estimator can be bounded.
To further account for the stochasticity of the process, a lower bound of the estimation is obtained using a concentration inequality. 
This involves an upper bound of the variance of the previous estimator. 
Interestingly, controlling for the variance of the estimator indirectly constrains the non-stationarity of the policy.
This prevents overfitting past non-stationary dynamics and better generalisation on future ones.
The algorithm in practice is implemented as a hyper-policy taking time as input and outputting the policy to be queried at each time.
Experiments show how this algorithm learns about the non-stationarity of the process to adapt to it, including when the non-stationarity is caused by delays.

\subsubsection{Extensive Empirical Evaluation}
Our last contribution is an extensive empirical evaluation of the algorithms proposed in this dissertation, as well as algorithms from the literature.
We design a wide range of test tasks, from balancing a pendulum to trading currency rates.
In these tasks, we vary the delay in length and nature to study the robustness of the algorithms. 
We consider constant, non-integer, stochastic, and multiple action delays. 
These empirical evaluations are useful for understanding the peculiarities of the algorithms and their preferred settings.


\subsubsection{Reference to Publications}
The aforementioned contributions can be found in the following papers. 
In \cite{liotet2021learning}, we present the model-based approach and analyse how model-based approaches may be too restrictive. 
For this work, I designed the approach, performed the theoretical analysis, and conducted most of the experimental analysis.
The imitation learning approach and analysis of the delay in smooth \glspl{mdp} can be found in \cite{liotet2022delayed}.
I have done part of the theoretical analysis for constant delays and done the extension to non-integer and stochastic delays. I have also conducted most of the experimental analysis.
The non-stationary memoryless approach comes from \cite{liotet2022lifelong}.
In this work, I conducted most of the theoretical analysis, part of the design of the approach, and all the empirical study. 
Finally, the results concerning the extension to multiple action delay and the analysis of the impact of changing the delay's distribution on the performance of the agent are the ground for a paper to be submitted. 
For this work, I have done all the theoretical and empirical analyses.





